\chapter{Le strutture}
\label{chap:strutture}

Gli array raccolgono molti dati \textbf{dello stesso tipo}. Spesso, però,
un'informazione è fatta di parti di tipo diverso: una carta da gioco ha un
valore (un numero) e un seme (un carattere); uno studente ha un nome (una
stringa), una matricola (un intero) e un voto (un carattere). Le
\textbf{strutture} (\code{struct}) permettono di aggregare variabili non
omogenee in un'unica struttura dati, da usare come un tutto. Sono anche il
mattone delle liste collegate del capitolo~\ref{chap:liste}.

\section{Definire una struttura}

\subsection{Il tipo \code{struct carta}}

Una carta da gioco italiana è identificata da un valore numerico, da 1 a 10,
e da un seme: oro, bastoni, coppe o spade. Possiamo dichiarare il tipo

\begin{lstlisting}[style=c]
struct carta {
    int  valore;          /* da 1 a 10 */
    char seme;            /* 'o', 'b', 'c', 's' */
};
\end{lstlisting}

\begin{itemize}
  \item \code{struct} è una parola chiave;
  \item \code{carta} è il nome assegnato alla struttura, detto
        \textbf{tag} (etichetta);
  \item \code{valore} e \code{seme} sono i \textbf{campi} (o membri) della
        struttura, ciascuno con il suo tipo.
\end{itemize}

Questa dichiarazione definisce il tipo di dato derivato \code{struct carta}
(capitolo~\ref{chap:tipi}): è un modello, e \textbf{nessuna memoria} viene
allocata in questo momento. La memoria viene allocata dalla dichiarazione
delle variabili:

\begin{lstlisting}[style=c]
struct carta c1, c2;     /* due variabili di tipo struct carta */
\end{lstlisting}

(Nelle slide, a questo punto, il tipo è chiamato \code{struct card}: è
rimasto il nome dell'esempio originale in inglese, ma il tipo dichiarato è
\code{struct carta}.)

\subsection{L'operatore punto}

Ai campi di una variabile struttura si accede con l'operatore \code{.}
(punto): \code{c1.valore} è il campo \code{valore} di \code{c1}, e si usa
come qualsiasi variabile \code{int}. Le strutture si possono anche
\textbf{assegnare} in blocco: \code{c2 = c1;} copia tutti i campi.

\begin{affianco}[0.5][c]
\begin{lstlisting}[style=c, name=carta.c]
#include <stdio.h>

typedef struct carta {
   int  valore;
   char seme;
} carta;

int main(void)
{
   carta c1, c2;

   c1.valore = 7;
   c1.seme = 'c';
   c2 = c1;                  /* copia tutti i campi */
   c2.valore = 1;
   printf("c1: %d di %c\n", c1.valore, c1.seme);
   printf("c2: %d di %c\n", c2.valore, c2.seme);
   printf("sizeof(carta) = %zu\n", sizeof(carta));
   return 0;
}
\end{lstlisting}
\risultato
\begin{console}
$ ./carta
c1: 7 di c
c2: 1 di c
sizeof(carta) = 8
\end{console}
\wtnota{Dopo la copia \code{c1} e \code{c2} sono indipendenti: cambiare
\code{c2.valore} non tocca \code{c1}.}
\end{affianco}

Il programma usa il \code{typedef} che vedremo fra poco. Ma perché la
struttura occupa 8 byte, se un \code{int} ne occupa 4 e un \code{char} 1? Il
processore accede più facilmente agli \code{int} che stanno a indirizzi
multipli di 4 (su alcuni processori è addirittura obbligatorio), e il
compilatore rispetta questa regola anche dentro le strutture. Pensiamo a un
array di carte: se ogni carta occupasse 5 byte, la seconda inizierebbe al
byte 5 e il suo \code{valore} non starebbe a un multiplo di 4. Per questo il
compilatore aggiunge alla fine della struttura 3 byte di
\textbf{riempimento} (\textit{padding}): ogni carta occupa 8 byte, e ogni
\code{valore} cade al posto giusto (figura~\ref{fig:carta}). Ecco perché la
dimensione di una struttura va sempre chiesta a \code{sizeof}, senza sommare
a mano quelle dei campi.

\begin{figure}[H]
  \centering
  \begin{tikzpicture}[font=\Lato\scriptsize, x=0.78cm, y=1cm]
    \foreach \base/\val/\sem in {0/7/c, 8/1/s} {
      \draw[draw=primary!65!black, fill=primary!3!white, line width=0.7pt]
        (\base,0) rectangle ++(4,0.55);
      \foreach \j in {1,2,3}
        \draw[draw=primary!35!white, densely dashed, line width=0.4pt]
          (\base+\j,0) -- ++(0,0.55);
      \node[font=\FiraCodeNL\footnotesize] at (\base+2,0.275) {\val};
      \draw[draw=accent!65!black, fill=accent!12!white, line width=0.7pt]
        (\base+4,0) rectangle ++(1,0.55);
      \node[font=\FiraCodeNL\scriptsize] at (\base+4.5,0.275) {'\sem'};
      \draw[draw=neutral!45!white, fill=neutral!12!white, line width=0.7pt]
        (\base+5,0) rectangle ++(3,0.55);
    }
    \foreach \k in {0,...,15}
      \node[mindir, anchor=north] at (\k+0.5,-0.02) {\k};
    % che cosa contiene la prima carta
    \draw[wtbrace] (4,0.68) -- (0,0.68);
    \node[mnota, anchor=south, align=center] at (2,0.92) {\code{valore}\\4 byte};
    \draw[draw=neutral!40!black, line width=0.6pt] (4.5,0.6) -- (4.5,0.9);
    \node[mnota, anchor=south, align=center] at (4.5,0.92) {\code{seme}\\1 byte};
    \draw[wtbrace] (8,0.68) -- (5,0.68);
    \node[mnota, anchor=south, align=center] at (6.5,0.92) {riempimento\\3 byte};
    \draw[wtbrace] (12,0.68) -- (8,0.68);
    \node[mnota, anchor=south, align=center, text=accent!70!black] at (10,0.92)
      {\code{mazzo[1].valore}\\inizia al byte 8};
    % le due carte
    \draw[wtbrace] (0,-0.42) -- (8,-0.42);
    \node[mnota, anchor=north] at (4,-0.66) {\code{mazzo[0]}: byte da 0 a 7};
    \draw[wtbrace] (8,-0.42) -- (16,-0.42);
    \node[mnota, anchor=north] at (12,-0.66) {\code{mazzo[1]}: byte da 8 a 15};
  \end{tikzpicture}
  \caption{Un array \code{carta mazzo[2]}: ogni carta occupa 8 byte (4 per
           \code{valore}, 1 per \code{seme} e 3 di riempimento), così il
           \code{valore} di ogni carta inizia a un multiplo di 4.}
  \label{fig:carta}
\end{figure}

\subsection{\code{typedef} e strutture}

Scrivere ogni volta \code{struct carta} è scomodo. Si può usare il
\code{typedef} (capitolo~\ref{chap:tipi}) per dare al tipo un nome più
breve, separatamente o direttamente nella definizione:

\begin{affianco}[0.5][c]
\begin{lstlisting}[style=c]
struct carta {
    int  valore;
    char seme;
};

typedef struct carta carta;
\end{lstlisting}
\risultato
\begin{lstlisting}[style=c]
typedef struct carta {
    int  valore;
    char seme;
} carta;
\end{lstlisting}
\end{affianco}

In entrambi i casi \code{carta} diventa sinonimo di \code{struct carta}, e le
variabili si dichiarano con \code{carta c1, c2;}. Il tag e il nome del
\code{typedef} possono coincidere, perché il C li tiene in elenchi separati.

\subsection{I nomi dei campi}

All'interno di una struttura i campi devono avere nomi diversi; campi di
strutture diverse, invece, possono avere lo stesso nome, perché si
distinguono dalla variabile a cui appartengono:

\begin{affianco}[0.5][c]
\begin{lstlisting}[style=c]
struct fruit {
   char   *name;
   int    calories;
};

struct vegetable {
   char    *name;
   int     calories;
};
\end{lstlisting}
\risultato
\begin{lstlisting}[style=c]
struct fruit      a;
struct vegetable  b;

a.calories = 38;
b.calories = 57;
\end{lstlisting}
\end{affianco}

\subsection{Strutture con e senza tag}

Il tag si può omettere, dichiarando subito le variabili:

\begin{lstlisting}[style=c]
struct {
   int    day, month, year;
   char   day_name[4];
   char   month_name[4];
} yesterday, today, tomorrow;
\end{lstlisting}

Senza tag, però, non si possono dichiarare altre variabili di quel tipo oltre
a \code{yesterday}, \code{today} e \code{tomorrow}: il tipo non ha un nome.
Con un tag, \code{struct date \{ \dots \};}, si possono dichiarare
liberamente variabili \code{struct date}. E con il \code{typedef} il tag non
serve: il nome lo dà il \code{typedef}.

\begin{affianco}[0.55][c]
\begin{lstlisting}[style=c, name=data.c]
#include <stdio.h>

typedef struct {
   int    day, month, year;
   char   day_name[4];
   char   month_name[4];
} date;

int main(void)
{
   date d1 = {31, 10, 2026, "Sat", "Oct"};

   printf("%s %d %s %d\n", d1.day_name, d1.day,
          d1.month_name, d1.year);
   return 0;
}
\end{lstlisting}
\risultato
\begin{console}
$ ./data
Sat 31 Oct 2026
\end{console}
\wtnota{Una struttura si inizializza come un array: i valori fra graffe, nello
stesso ordine dei campi.}
\end{affianco}

\subsection{Che cosa non si può fare}

Le strutture si possono assegnare, passare alle funzioni e restituire, ma non
si possono \textbf{confrontare} con \code{==}, né stampare con un'unica
conversione di \code{printf}. Proviamo a confrontare due carte:

\begin{affianco}[0.5][c]
\begin{lstlisting}[style=c, name=confronta.c]
typedef struct carta {
   int  valore;
   char seme;
} carta;

int main(void)
{
   carta c1 = {7, 'c'}, c2 = {7, 'c'};

   if (c1 == c2)
      return 1;
   return 0;
}
\end{lstlisting}
\risultato
Il compilatore non sa che cosa significhi \code{==} fra due strutture, e
rifiuta il programma (qui sotto). Il confronto va scritto campo per campo:

\vspace{4pt}
\begin{lstlisting}[style=c]
if (c1.valore == c2.valore
        && c1.seme == c2.seme)
   return 1;
\end{lstlisting}
\end{affianco}

\begin{console}
$ gcc -Wall -o confronta confronta.c
confronta.c: In function (*@‘@*)main(*@’@*):
confronta.c:10:11: error: invalid operands to binary == (have (*@‘@*)carta(*@’@*) and (*@‘@*)carta(*@’@*))
   10 |    if (c1 == c2)
      |           ^~
\end{console}

\section{Array di strutture}

\subsection{Un tipo come gli altri}

Una volta definito, un tipo struttura si usa come qualsiasi altro tipo: si
possono dichiarare array di strutture, puntatori a strutture, e allocarle
dinamicamente.

\begin{lstlisting}[style=c]
typedef struct {
   float   re;
   float   im;
} complex;

complex a, b, c[100], *d;
...
d = (complex *) malloc(n * sizeof(complex));
\end{lstlisting}

\code{c[100]} è un array di 100 numeri complessi: \code{c[i].re} è la parte
reale dell'elemento \code{i}. \code{d} punta a un array di \code{n} complessi
allocato nello heap, e anche lì si scrive \code{d[i].re}.

\subsection{Il programma \code{class\_info}}

L'esempio delle slide gestisce le informazioni su una classe. In un header,
incluso da ogni modulo del programma, si definiscono la dimensione massima e
il tipo \code{student}; la funzione \code{fail} conta gli studenti con il
voto \code{'F'} (insufficiente, nel sistema americano):

\begin{affianco}[0.5][t]
\begin{lstlisting}[style=h, name=class_info.h]
#define   CLASS_SIZE    100

typedef struct {
   char     *last_name;
   int      student_id;
   char     grade;
} student;

int fail(student *class, int N);
\end{lstlisting}
\begin{lstlisting}[style=c, name=fail.c]
/* Count the failing grades. */
#include "class_info.h"

int fail(student *class, int N)
{
   int   i, cnt = 0;

   for (i = 0; i < N; ++i)
      cnt += class[i].grade == 'F';
   return cnt;
}
\end{lstlisting}
\risultato
\begin{lstlisting}[style=c, name=main.c]
#include <stdio.h>
#include <stdlib.h>
#include "class_info.h"

int main(void)
{
   int       cnt, i;
   student   tmp, *class;

   class = malloc(CLASS_SIZE * sizeof(student));
   tmp.grade = 'A';
   tmp.last_name = "Casanova";
   tmp.student_id = 910017;
   class[0] = tmp;
   tmp.grade = 'F';
   tmp.last_name = "Superman";
   tmp.student_id = 910018;
   class[1] = tmp;
   tmp.grade = 'C';
   tmp.last_name = "Wonderwoman";
   tmp.student_id = 910019;
   class[2] = tmp;
   /* Enter nothing for the remainder of the class. */
   for (i = 3; i < CLASS_SIZE; ++i) {
      class[i].grade = ' ';
      class[i].last_name = "";
      class[i].student_id = 0;
   }
   cnt = fail(class, CLASS_SIZE);
   printf("Number of failing grades: %d\n", cnt);
   free(class);
   return 0;
}
\end{lstlisting}
\end{affianco}

\begin{affianco}[0.5][c]
\begin{console}
$ gcc -Wall -o class_info main.c fail.c
$ ./class_info
Number of failing grades: 1
\end{console}
\risultato
Nelle slide ci sono due versioni di \code{main}: una con l'array
\code{student class[CLASS\_SIZE]} dichiarato normalmente, e questa, con
l'array allocato dinamicamente. \code{fail} va bene per entrambe, perché
riceve un puntatore (capitolo~\ref{chap:puntatori}).
\end{affianco}

\subsection{Come funziona}

\begin{itemize}
  \item \code{tmp} è una variabile di appoggio: se ne riempiono i campi e
        la si copia in \code{class[0]}, poi la si riempie di nuovo e la si
        copia in \code{class[1]}, e così via (figura~\ref{fig:classe}).
  \item La copia \code{class[0] = tmp} copia \textbf{i campi}: per
        \code{last\_name}, che è un puntatore, copia l'indirizzo, non la
        stringa. Qui va bene, perché ogni studente punta a una stringa
        costante diversa.
  \item Gli altri 97 studenti hanno voto \code{' '}, cognome \code{""} (la
        stringa vuota) e matricola 0.
  \item In \code{fail}, \code{class[i].grade == 'F'} vale 1 o 0, e
        \code{cnt += \dots} lo somma: si contano i voti \code{'F'} senza
        bisogno di un \code{if}.
\end{itemize}

\begin{figure}[H]
  \centering
  \begin{tikzpicture}[font=\Lato\scriptsize]
    \mchiave{2.0cm}
    \oggetto[anchor=north west]{s0}{(0,0)}{3.3cm}{class[0]}{last\_name/{}, student\_id/910017, grade/{\code{'A'}}}
    \oggetto[anchor=north west]{s1}{([yshift=0.7pt]s0-box.south west)}{3.3cm}{class[1]}{last\_name/{}, student\_id/910018, grade/{\code{'F'}}}
    \oggetto[anchor=north west]{s2}{([yshift=0.7pt]s1-box.south west)}{3.3cm}{class[2]}{last\_name/{}, student\_id/910019, grade/{\code{'C'}}}
    \node[mnota, anchor=north, align=center] at ([yshift=-0.05cm]s2-box.south) {$\vdots$\\fino a \code{class[99]}};
    % le stringhe a cui puntano i campi last_name
    \node[mvar, anchor=west, minimum width=2.5cm] (t0) at ([xshift=1.3cm]s0-1.east) {"Casanova"};
    \node[mvar, anchor=west, minimum width=2.5cm] (t1) at ([xshift=1.3cm]s1-1.east) {"Superman"};
    \node[mvar, anchor=west, minimum width=2.5cm] (t2) at ([xshift=1.3cm]s2-1.east) {"Wonderwoman"};
    \foreach \ist in {0,1,2} {
      \draw[mrif] ([xshift=-0.45cm]s\ist-1.east) -- (t\ist.west);
      \node[mpunto] at ([xshift=-0.45cm]s\ist-1.east) {};
    }
    \node[mnota, anchor=south] at ([yshift=2pt]t0.north) {stringhe costanti};
    % il puntatore class
    \node[mvar, minimum width=0.7cm] (pc) at ([xshift=-1.3cm]s0.west) {};
    \node[mnome, anchor=east] at ([xshift=-2pt]pc.west) {class};
    \riferimento{pc}{-- (s0.west)}
  \end{tikzpicture}
  \caption{L'array \code{class} nello heap: 100 strutture \code{student}
           consecutive, di 16 byte ciascuna, e le prime tre contengono i dati
           degli studenti. Il campo \code{last\_name} contiene solo
           l'indirizzo della stringa.}
  \label{fig:classe}
\end{figure}

\section{L'operatore freccia}

\subsection{\code{p->campo}}

Quando si ha un \textbf{puntatore} a una struttura, per accedere a un campo
si dovrebbe scrivere \code{(*p).grade}: prima si dereferenzia il puntatore,
poi si sceglie il campo. Le parentesi sono necessarie, perché \code{.} ha la
precedenza su \code{*}: \code{*p.grade} vorrebbe dire \code{*(p.grade)}. Per
questo il C ha un operatore apposito, la \textbf{freccia} \code{->}:

\[
  \code{p->grade} \quad\text{equivale a}\quad \code{(*p).grade}
\]

Le slide riscrivono \code{fail} in altri due modi, entrambi con la freccia:

\begin{affianco}[0.5][c]
\begin{lstlisting}[style=c, name=fail2.c]
int fail(student *class, int N)
{
   int   i, cnt = 0;

   for (i = 0; i < N; ++i)
      cnt += (class + i)->grade == 'F';
   return cnt;
}
\end{lstlisting}
\risultato
\begin{lstlisting}[style=c, name=fail3.c]
int fail(student *class, int N)
{
   int       i, cnt = 0;
   student  *tmp = class;

   for (i = 0; i < N; ++i, tmp++)
      cnt += tmp->grade == 'F';
   return cnt;
}
\end{lstlisting}
\end{affianco}

Nella prima \code{class + i} è l'indirizzo dell'elemento \code{i} (per
l'aritmetica dei puntatori avanza di una \textbf{struttura} intera alla
volta), e \code{->grade} ne prende il voto. Nella seconda un puntatore
\code{tmp} scorre l'array: \code{tmp++} lo porta sullo studente successivo.
Tutte e tre le versioni danno lo stesso risultato (verificato).

\subsection{Espressioni con i puntatori a strutture}

Le slide riportano una tabella del libro di Kelley e Pohl, nel quale il tipo
è dichiarato con un tag, \code{struct student \{ \dots \};}, e ha gli stessi
campi di \code{class\_info.h}. Con le dichiarazioni e le assegnazioni

\begin{lstlisting}[style=c]
struct student  tmp, *p = &tmp;

tmp.grade = 'A';
tmp.last_name = "Casanova";
tmp.student_id = 910017;
\end{lstlisting}

valgono queste equivalenze (figura~\ref{fig:freccia}):

\begin{center}
\renewcommand{\arraystretch}{1.15}
\begin{tabular}{@{}>{\FiraCodeNL}l>{\FiraCodeNL}l l@{}}
  \toprule
  \normalfont\textbf{espressione} & \normalfont\textbf{equivalente} & \textbf{valore} \\
  \midrule
  tmp.grade               & p->grade               & \code{A} \\
  tmp.last\_name          & p->last\_name          & \code{Casanova} \\
  (*p).student\_id        & p->student\_id         & 910017 \\
  *p->last\_name + 1      & (*(p->last\_name)) + 1 & \code{D} \\
  *(p->last\_name + 2)    & (p->last\_name)[2]     & \code{s} \\
  \bottomrule
\end{tabular}
\end{center}

Le ultime due sono le più istruttive. In \code{*p->last\_name + 1} la freccia
lega più di tutti: si prende \code{p->last\_name}, la stringa; \code{*} ne
dà il primo carattere, \code{'C'}; e \code{+ 1} dà il carattere successivo,
\code{'D'}. In \code{*(p->last\_name + 2)}, invece, si avanza di due caratteri
nella stringa e si prende quello che c'è: la \code{s} di Ca\textbf{s}anova.

\begin{figure}[H]
  \centering
  \begin{tikzpicture}[font=\Lato\scriptsize]
    \mchiave{2.3cm}
    \oggetto{t}{(3.0,0)}{3.6cm}{tmp}{last\_name/{}, student\_id/910017, grade/{\code{'A'}}}
    \node[mnome, anchor=east] at (-0.15,-0.6) {p};
    \node[mvar, minimum width=0.7cm] (pp) at (0.3,-0.6) {};
    \riferimento{pp}{-- (t-box.west)}
    \noindici
    \celle[mcella, minimum width=0.45cm, inner xsep=0pt]{nm}{(7.0,-0.1)}{C,a,s,a,n,o,v,a,\textbackslash 0}
    \draw[mrif] ([xshift=-0.45cm]t-1.east) to[out=0, in=180] (nm-0.west);
    \node[mpunto] at ([xshift=-0.45cm]t-1.east) {};
    \draw[wtarrow] ([yshift=-0.6cm]nm-0.south) -- ([yshift=-1pt]nm-0.south);
    \node[mnota, anchor=north west] at ([xshift=-0.3cm, yshift=-0.6cm]nm-0.south)
      {\code{*p->last\_name} è la \code{C}; \code{+ 1} dà la \code{D}};
    \draw[wtarrowacc] ([yshift=0.55cm]nm-2.north) -- ([yshift=1pt]nm-2.north);
    \node[mnota, anchor=south, text=accent!70!black] at ([yshift=0.55cm]nm-2.north)
      {\code{p->last\_name + 2}: la \code{s}};
  \end{tikzpicture}
  \caption{\code{p} punta alla struttura \code{tmp}; il campo
           \code{last\_name} punta a sua volta alla stringa.}
  \label{fig:freccia}
\end{figure}

\subsection{Una verifica}

Il programma seguente dichiara il tipo con il tag, come nel libro, e stampa i
valori delle espressioni della tabella:

\begin{affianco}[0.6][c]
\begin{lstlisting}[style=c, name=puntatori.c]
#include <stdio.h>

struct student {
   char *last_name;
   int   student_id;
   char  grade;
};

int main(void)
{
   struct student tmp, *p = &tmp;

   tmp.grade = 'A';
   tmp.last_name = "Casanova";
   tmp.student_id = 910017;
   printf("%c\n", p -> grade);
   printf("%s\n", p -> last_name);
   printf("%d\n", (*p).student_id);
   printf("%c\n", *p -> last_name + 1);
   printf("%c\n", *(p -> last_name + 2));
   return 0;
}
\end{lstlisting}
\risultato
\begin{console}
$ ./puntatori
A
Casanova
910017
D
s
\end{console}
\wtnota{Come nella tabella delle slide, attorno a \code{->} ci sono degli
spazi: sono permessi, perché \code{->} è un unico token, ma di solito si
scrive \code{p->grade}.}
\end{affianco}

\code{*p -> last\_name + 1} è un numero, il codice del carattere
\code{'C'} più 1, cioè quello di \code{'D'}: per vederlo come carattere si
stampa con \code{\%c}.

\section{Strutture di strutture}

\subsection{Campi che sono strutture}

Un campo può essere a sua volta una struttura, o un puntatore a una
struttura. L'esempio delle slide descrive un dipendente:

\begin{lstlisting}[style=c]
typedef struct {
   char           name[20];
   int            number;
} department;

typedef struct {
   char           street_name[30];
   int            no;
   char           postcode[10];
} address;

typedef struct {
   char           name[25];
   int            employee_id;
   department     dept;          /* una struttura dentro la struttura */
   address        *home;         /* un puntatore a un'altra struttura */
   double         salary;
} employee_data;
\end{lstlisting}

(Nelle slide \code{name} di \code{department} e \code{street\_name} di
\code{address} sono dichiarati \code{char}, cioè un solo carattere: per
contenere un nome servono degli array, come qui.) Il reparto è contenuto
nella struttura del dipendente; l'indirizzo di casa, invece, sta altrove, e
il dipendente ne contiene solo il puntatore (figura~\ref{fig:annidate}).
Per arrivare a un campo interno si concatenano gli operatori, da sinistra a
destra: \code{e.dept.number} equivale a \code{(e.dept).number}, e con un
puntatore \code{p->dept.number} equivale a \code{(p->dept).number};
\code{e.home->no} è il numero civico.

\begin{figure}[H]
  \centering
  \begin{tikzpicture}[font=\Lato\scriptsize,
      riga/.style={mriga, minimum width=5.9cm,
        text width=\dimexpr5.9cm-6pt\relax, anchor=north}]
    % la struttura e, con la struttura dept al suo interno
    \node[mtesta, minimum width=5.9cm, anchor=north west] (e) at (0,0) {employee\_data e};
    \node[riga] (e-1) at ([yshift=0.7pt]e.south)
      {\makebox[2.1cm][l]{\color{primary!55!black}name}\code{"Mario Rossi"}};
    \node[riga] (e-2) at ([yshift=0.7pt]e-1.south)
      {\makebox[2.1cm][l]{\color{primary!55!black}employee\_id}1234};
    \node[riga, minimum height=1.85cm] (e-3) at ([yshift=0.7pt]e-2.south) {};
    \node[anchor=north west, font=\FiraCodeNL\scriptsize, text=primary!55!black,
          inner sep=3pt] at ([yshift=-2pt]e-3.north west) {dept};
    \mchiave{1.2cm}
    \oggetto[anchor=north west, fill=accent!12!white]{d}{([xshift=2.2cm, yshift=-0.2cm]e-3.north west)}{3.2cm}{department}{name/{\code{"Vendite"}}, number/1}
    \node[riga] (e-4) at ([yshift=0.7pt]e-3.south)
      {\makebox[2.1cm][l]{\color{primary!55!black}home}};
    \node[riga] (e-5) at ([yshift=0.7pt]e-4.south)
      {\makebox[2.1cm][l]{\color{primary!55!black}salary}1850.0};
    % la struttura address, altrove in memoria
    \mchiave{2.0cm}
    \oggetto[anchor=west]{a}{([xshift=1.6cm]e-4.east)}{4.1cm}{address}{street\_name/{\code{"Via Roma"}}, no/12, postcode/{\code{"80100"}}}
    \draw[mrif] ([xshift=-0.5cm]e-4.east) -- (a.west);
    \node[mpunto] at ([xshift=-0.5cm]e-4.east) {};
    % come si arriva a due campi interni
    \node[mnota, anchor=west, text=accent!70!black] (l1) at ([xshift=0.6cm]e-3.east |- d-2) {\code{e.dept.number}};
    \draw[wtarrowacc] (l1.west) -- (d-2.east);
    \node[mnota, anchor=west, text=accent!70!black] (l2) at ([xshift=0.6cm]a-2.east) {\code{e.home->no}};
    \draw[wtarrowacc] (l2.west) -- (a-2.east);
  \end{tikzpicture}
  \caption{\code{dept} è una struttura dentro \code{e}; \code{home} punta a una
           struttura che sta altrove. Le etichette mostrano le espressioni che
           arrivano a due campi interni.}
  \label{fig:annidate}
\end{figure}

\section{Strutture e funzioni}

\subsection{Due modi di passare una struttura}

Supponiamo di voler scrivere una funzione che aggiorna le informazioni di un
dipendente, per esempio il numero del reparto. Si può procedere in due
maniere:

\begin{affianco}[0.5][c]
\begin{lstlisting}[style=c]
employee_data update(employee_data e)
{
   ...
   e.dept.number = n;
   ...
   return e;
}

/* uso: */
e = update(e);
\end{lstlisting}
\risultato
\begin{lstlisting}[style=c]
void update(employee_data *p)
{
   ...
   p->dept.number = n;
   ...
}

/* uso: */
update(&e);
\end{lstlisting}
\end{affianco}

\begin{itemize}
  \item Nel primo metodo \code{e} è passata \textbf{per valore}: la funzione
        riceve una copia di tutta la struttura, la modifica e la restituisce,
        e il chiamante deve assegnare il risultato. Se la struttura è grande,
        creare la copia (due volte: all'andata e al ritorno) può essere
        oneroso.
  \item Nel secondo metodo si passa l'\textbf{indirizzo} di \code{e}: non si
        crea nessuna copia, e la funzione modifica direttamente la struttura
        del chiamante. È il metodo in generale più efficiente.
\end{itemize}

\subsection{Una prova}

Il programma seguente completa le due funzioni (con il nuovo numero di
reparto come parametro \code{n}) e mostra la differenza. In cima al file ci
sono i tre \code{typedef} della sezione precedente, che qui non sono
ripetuti:

\begin{affianco}[0.58][c]
\codicepiccolo
\begin{lstlisting}[style=c, name=impiegati.c]
#include <stdio.h>

/* ... i typedef di department, address ed employee_data ... */

employee_data update(employee_data e, int n)    /* per valore */
{
   e.dept.number = n;
   return e;
}

void update_p(employee_data *p, int n)          /* per indirizzo */
{
   p->dept.number = n;
}

int main(void)
{
   employee_data e = {"Mario Rossi", 1234, {"Vendite", 1}, NULL, 1850.0};

   update(e, 2);                    /* il risultato non viene usato */
   printf("dopo update(e, 2):       reparto %d\n", e.dept.number);
   e = update(e, 3);
   printf("dopo e = update(e, 3):   reparto %d\n", e.dept.number);
   update_p(&e, 4);
   printf("dopo update_p(&e, 4):    reparto %d\n", e.dept.number);
   printf("sizeof(employee_data) = %zu byte\n", sizeof(employee_data));
   return 0;
}
\end{lstlisting}
\risultato
\begin{console}
$ ./impiegati
dopo update(e, 2):       reparto 1
dopo e = update(e, 3):   reparto 3
dopo update_p(&e, 4):    reparto 4
sizeof(employee_data) = 72 byte
\end{console}
\end{affianco}

\begin{itemize}
  \item \code{update(e, 2)} da sola non cambia nulla: modifica la copia, che
        poi viene buttata via perché il risultato non è assegnato.
  \item \code{e = update(e, 3)} funziona, ma copia 72 byte all'andata e 72
        al ritorno.
  \item \code{update\_p(\&e, 4)} copia solo un indirizzo, 8 byte, e modifica
        \code{e} direttamente (figura~\ref{fig:struct-funzioni}).
  \item La struttura si inizializza con le graffe, anche quella interna
        (\code{\{"Vendite", 1\}}); \code{home} vale \code{NULL}, perché in
        questa prova l'indirizzo non c'è.
\end{itemize}

\subsection{Copia o indirizzo}

La figura~\ref{fig:struct-funzioni} mette a confronto i due modi. Per valore,
\code{update} lavora su una copia di tutti i 72 byte della struttura, e il
risultato va ricopiato in \code{e}; per indirizzo, \code{update\_p} riceve
soltanto gli 8 byte di un puntatore e, attraverso di esso, modifica proprio
la \code{e} di \code{main}.

\begin{figure}[H]
  \centering
  \begin{tikzpicture}[font=\Lato\scriptsize]
    \node[font=\Lato\scriptsize\bfseries] at (2.2,1.1) {\code{e = update(e, 3)}: per valore};
    \node[wtnodeplain, minimum width=2.0cm, minimum height=1.0cm] (m1) at (0.9,0) {\code{e}\\72 byte};
    \node[wtnodeplain, minimum width=2.0cm, minimum height=1.0cm, fill=accent!10!white] (f1) at (3.9,0) {copia di \code{e}\\72 byte};
    \draw[mcerca] ([yshift=0.15cm]m1.east) -- node[above, mnota] {copia} ([yshift=0.15cm]f1.west);
    \draw[mcerca] ([yshift=-0.15cm]f1.west) -- node[below, mnota] {copia} ([yshift=-0.15cm]m1.east);
    \node[mnota] at (0.9,-0.8) {\code{main}};
    \node[mnota] at (3.9,-0.8) {\code{update}};
    \begin{scope}[xshift=7.6cm]
      \node[font=\Lato\scriptsize\bfseries] at (2.2,1.1) {\code{update\_p(\&e, 4)}: per indirizzo};
      \node[wtnodeplain, minimum width=2.0cm, minimum height=1.0cm] (m2) at (0.9,0) {\code{e}\\72 byte};
      \node[mvar, fill=accent!10!white] (p2) at (3.9,0) {};
      \node[mnome, anchor=west] at ([xshift=2pt]p2.east) {p};
      \riferimento{p2}{-- (m2.east)}
      \node[mnota] at (0.9,-0.8) {\code{main}};
      \node[mnota] at (3.9,-0.8) {\code{update\_p}: 8 byte};
    \end{scope}
  \end{tikzpicture}
  \caption{Passare una struttura per valore copia tutti i suoi byte;
           passarla per indirizzo copia solo un puntatore.}
  \label{fig:struct-funzioni}
\end{figure}

Quando una funzione deve solo \textbf{leggere} una struttura grande, si passa
comunque l'indirizzo, ma dichiarando il parametro \code{const}:

\begin{lstlisting}[style=c]
void stampa(const employee_data *p);    /* legge *p, ma non può modificarlo */
\end{lstlisting}

Così si evita la copia, e il compilatore garantisce che la funzione non
modifichi la struttura (capitolo~\ref{chap:stringhe}).

\insintesi{
  Una \textbf{struttura} aggrega campi di tipo diverso: la dichiarazione
  \code{struct tag \{ \dots \};} definisce un tipo (senza allocare memoria),
  le variabili si dichiarano con \code{struct tag v;} o, con un
  \code{typedef}, con il solo nome del tipo. Ai campi si accede con il
  \textbf{punto} (\code{v.campo}) e, attraverso un puntatore, con la
  \textbf{freccia} (\code{p->campo}, cioè \code{(*p).campo}). Le strutture si
  inizializzano con le graffe, si assegnano in blocco (copiando i campi, e
  per i puntatori solo l'indirizzo) ma non si confrontano con \code{==}; la
  loro dimensione, che può includere byte di riempimento, va chiesta a
  \code{sizeof}. Si possono fare array di strutture, anche dinamici, e
  strutture che contengono strutture o puntatori a strutture. Alle funzioni
  conviene passare l'indirizzo di una struttura, invece di una copia.}

\section{Riferimenti}

\begin{itemize}
  \item \textbf{A. Kelley, I. Pohl}, \textit{A Book on C}: capitolo~9
        (strutture e unioni).
  \item \textbf{Struct declaration} --- cppreference.com. \\
        \urlx{https://en.cppreference.com/w/c/language/struct}
\end{itemize}
